% Part: incompleteness
% Chapter: introduction
% Section: definitions

\documentclass[../../../include/open-logic-section]{subfiles}

\begin{document}

% SEGMENT OLP-0277-S01

\olfileid{inc}{int}{def}

\olsection{வரையறைகள்}

கணிதத்தை முறைப்படுத்தி, அந்த முறைப்படுத்தல் முரணற்றதாகவும் நிறைவானதாகவும்
உள்ளது என்று காட்டும் ஹில்பர்டின் திட்டத்தைச் செயல்படுத்துவதற்கு, முதலில் ஒரு
மொழி, தருக்கக் கட்டமைப்பு, அடிகோள் முறைமை ஆகியவற்றைத் தேர்ந்தெடுக்க வேண்டும்.
நமது நோக்கங்களுக்காக, கணிதத்தை ஒரு முதல்தர மொழியில் முறைப்படுத்தலாம் என்று
கொள்வோம். அதாவது, முதல்தர தருக்கத்தின் இணைப்பிகளுடனும் அளவையடைகளுடனும்
சேர்ந்து கணிதக் கூற்றுகளை வெளிப்படுத்த உதவும் !!{constant}s,
!!{function}s, !!{predicate}s ஆகியவற்றின் ஏதோவொரு கணம் உள்ளது என்க.
அத்தகைய மொழி உள்ளது என்று பெரும்பாலோர் ஒப்புக்கொள்கின்றனர்: $\in$ மட்டுமே
தருக்கமல்லாத குறியாக அமையும் கணக் கோட்பாட்டு மொழி. இவ்வளவு எளிய மொழி
இவ்வளவு வெளிப்படுத்தும் திறனுள்ளது என்ற கூற்று முதற்பார்வையில் மிகவும்
நம்பவியலாததாகத் தோன்றும். ஏறத்தாழ கணிதம் முழுவதையும் இந்த மிகக் குறுகிய
சொற்றொகுதியால் வெளிப்படுத்தலாம் என்பதை நிறுவப் பெருமளவு பணி தேவைப்பட்டது.
இப்போதைக்கு எளிமை கருதி, இயல் எண்கள்~$\Nat$ பற்றி மட்டும் பேசும் கணிதப்
பகுதியான எண்கணிதத்துடன் நமது விவாதத்தை வரம்பிடுவோம். எண்கணித உண்மைகளை
வெளிப்படுத்துவதற்கான இயல்பான மொழி~$\Lang L_A$. $\Lang L_A$ இல் ஒரே ஓர்
ஈரிட !!{predicate}~$<$, ஒரே ஒரு !!{constant}~$\Obj 0$, ஓரிட
!!{function}~$\prime$, ஈரிட !!{function}s~$+$,~$\times$ ஆகியவை உள்ளன.

\begin{defn}
!!{sentence}s இன் கணம்~$\Gamma$, பின்விளைவின்கீழ் மூடப்பட்டிருந்தால்,
அதாவது $\Gamma = \Setabs{!A}{\Gamma \Entails !A}$ எனில், அது ஒரு
\emph{கோட்பாடு} ஆகும்.
\end{defn}

கோட்பாடுகளைக் குறிப்பிடுவதற்கு இரண்டு எளிய வழிகள் உள்ளன. ஏதாவது ஒரு
!!{structure} இல் மெய்யாகும் !!{sentence}s இன் கணமாகக் குறிப்பிடுவது
ஒரு வழி. எடுத்துக்காட்டாக, !!{domain} ஆனது~$\Nat$ ஆகவும், தருக்கமல்லாத
எல்லாக் குறிகளும் எதிர்பார்த்தவாறே விளக்கப்படுவதாகவும் இருக்கும்
$\Lang L_A$ மொழிக்கான !!{structure} ஐக் கருதுக.

\begin{defn}\ollabel{def:standard-model}
\emph{எண்கணிதத்தின் தரநிலை மாதிரி} என்பது பின்வருமாறு வரையறுக்கப்படும்
!!{structure}~$\Struct N$ ஆகும்:
\begin{enumerate}
\item $\Domain N = \Nat$
\item $\Assign{\Obj 0}{N} = 0$
\item எல்லா $n \in \Nat$ க்கும் $\Assign{\Obj \prime}{N}(n) = n + 1$
\item எல்லா $n, m \in \Nat$ க்கும் $\Assign{\Obj +}{N}(n, m) = n + m$
\item எல்லா $n, m \in \Nat$ க்கும் $\Assign{\Obj \times}{N}(n, m) = n\cdot m$
\item $\Assign{\Obj <}{N} = \Setabs{\tuple{n, m}}{n \in \Nat, m \in
  \Nat, n < m}$
\end{enumerate}
\end{defn}

`$\times$' க்கும் `$\cdot$' க்கும் உள்ள வேறுபாட்டைக் கவனிக்கவும்:
$\times$ என்பது எண்கணித மொழியிலுள்ள ஒரு குறி. பெருக்கலை நினைவூட்டுவதற்காக
அதைத் தேர்ந்தெடுத்தோம்; ஆனால் $\times$ என்பது பெருக்கல் செயல் அல்ல, ஓர்
ஈரிடச் சார்புக் குறி (முறைப்படியான பெயரில், $\Obj f^2_1$). மாறாக,
$\cdot$ சாதாரணப் பெருக்கல் சார்பே \emph{ஆகும்}. $n \cdot m$ போன்ற ஒன்றைக்
காணும்போது, $n$,~$m$ எண்களின் பெருக்கற்பலனைக் குறிக்கிறோம்; $x \times y$
போன்ற ஒன்றைக் காணும்போது, எண்கணித மொழியிலுள்ள ஒரு சொல்லைப் பற்றிப் பேசுகிறோம்.
தரநிலை மாதிரியில், $\times$ என்ற சார்புக் குறி இயல் எண்கள் மீதான
$\cdot$ என்ற சார்பாக விளக்கப்படுகிறது. கூட்டலுக்கு, எண்கணித மொழியின்
சார்புக் குறியாகவும் இயல் எண்கள் மீதான கூட்டல் சார்பாகவும் $+$ ஐயே
பயன்படுத்துகிறோம். இங்கு எது குறிக்கப்படுகிறது என்பதைச் சூழலிலிருந்து
தீர்மானிக்க வேண்டும்.

\begin{defn}
\emph{மெய் எண்கணிதக் கோட்பாடு} என்பது எண்கணிதத்தின் தரநிலை மாதிரியில்
நிறைவுறுத்தப்படும் !!{sentence}s இன் கணம்; அதாவது,
\[
\Th{TA} = \Setabs{!A}{\Sat{N}{!A}}.
\]
\end{defn}

$\Th{TA}$ ஒரு கோட்பாடு. ஏனெனில் $\Th{TA} \Entails !A$ எனில்,
$\Th{TA}$ ஐ நிறைவுறுத்தும் ஒவ்வொரு !!{structure} இலும்~$!A$
நிறைவுறுத்தப்படுகிறது. $\Sat{N}{\Th{TA}}$ என்பதால் $\Sat{N}{!A}$;
ஆகவே $!A \in \Th{TA}$.

ஒரு கோட்பாடு~$\Gamma$ ஐக் குறிப்பிடுவதற்கான மற்றொரு வழி, ஏதாவது ஒரு
வாக்கியக் கணம்~$\Gamma_0$ இன் பின்விளைவுகளாக அமையும் !!{sentence}s இன்
கணமாகக் குறிப்பிடுவதாகும். அப்போது $\Gamma$ என்பது பின்விளைவின்கீழ்
$\Gamma_0$ இன் ``மூடல்.'' $\Gamma_0$ வெளிப்படையாகக் குறிப்பிடப்பட்டால்,
எடுத்துக்காட்டாக $\Gamma_0$ இன் !!{element}s பட்டியலிடப்பட்டால் மட்டுமே, இவ்வழியில்
ஒரு கோட்பாட்டைக் குறிப்பிடுவது பயனுடையது. குறைந்தபட்சம் $\Gamma_0$
தீர்மானிக்கத்தக்கதாக இருக்க வேண்டும்; அதாவது ஓர் !!a{sentence} ஆனது
$\Gamma_0$ இன் உறுப்பா இல்லையா என்பதைச் சோதிக்கும் கணிக்கத்தக்க சோதனை
இருக்க வேண்டும். $\Gamma_0$ இலுள்ள !!{sentence}s ஐ~$\Gamma$ வின்
\emph{அடிகோள்கள்} என்றும், $\Gamma$ ஆனது~$\Gamma_0$ ஆல்
\emph{அடிகோளாக்கப்படுகிறது} என்றும் கூறுகிறோம்.

\begin{defn}
ஒரு கோட்பாடு~$\Gamma$ ஆனது~$\Gamma_0$ ஆல் \emph{அடிகோளாக்கப்படுகிறது}
அப்போதும் அப்போது மட்டுமே
\[
\Gamma = \Setabs{!A}{\Gamma_0 \Entails !A}
\]
\end{defn}

\begin{defn}
பின்வரும் வாக்கியங்களால் அடிகோளாக்கப்படும் $\Th{Q}$ கோட்பாடு
``ராபின்சனின் $\Th{Q}$'' என அறியப்படுகிறது; இது மிகவும் எளிய எண்கணிதக்
கோட்பாடு.
\begin{align*}
& \lforall[x][\lforall[y][(\eq[x'][y'] \lif \eq[x][y])]] \tag{$!Q_1$}\\
& \lforall[x][\eq/[\Obj 0][x']] \tag{$!Q_2$}\\
& \lforall[x][(\eq[x][\Obj 0] \lor \lexists[y][\eq[x][y']])] \tag{$!Q_3$}\\
& \lforall[x][\eq[(x + \Obj 0)][x]] \tag{$!Q_4$}\\
& \lforall[x][\lforall[y][\eq[(x + y')][(x + y)']]] \tag{$!Q_5$}\\
& \lforall[x][\eq[(x \times \Obj 0)][\Obj 0]] \tag{$!Q_6$}\\
& \lforall[x][\lforall[y][\eq[(x \times y')][((x \times y) + x)]]] \tag{$!Q_7$}\\
& \lforall[x][\lforall[y][(x < y \liff \lexists[z][\eq[(z' + x)][y]])]] \tag{$!Q_8$}
\end{align*}
!!{sentence}s இன் கணம்~$\{!Q_1, \dots, !Q_8\}$ ஆனது~$\Th{Q}$ இன்
அடிகோள்களாகும்; ஆகவே அவற்றின் பின்விளைவுகளான எல்லா !!{sentence}s ஐயும்
$\Th{Q}$ கொண்டுள்ளது:
\[
\Th{Q} = \Setabs{!A}{\{!Q_1, \dots, !Q_8\} \Entails !A}.
\]
\end{defn}

\begin{defn}
$!A(x)$ ஆனது~$\Lang L_A$ இலுள்ள !!a{formula} என்றும், அதன் கட்டற்ற
மாறிகள்~$x$, $y_1$, \dots, $y_n$ என்றும் கொள்க. அப்போது பின்வரும்
வடிவிலுள்ள எந்த !!{sentence} உம் \emph{தொகுத்தறிதல் வார்ப்புருவின்}
ஒரு நிகழ்வாகும்:
\[
\lforall[y_1][\dots\lforall[y_n][((!A(\Obj 0) \land \lforall[x][(!A(x)
\lif !A(x'))]) \lif \lforall[x][!A(x)])]]
\]

\emph{பியானோ எண்கணிதம்}~$\Th{PA}$ என்பது $\Th{Q}$ இன் அடிகோள்களுடனும்
தொகுத்தறிதல் வார்ப்புருவின் எல்லா நிகழ்வுகளுடனும் அடிகோளாக்கப்படும் கோட்பாடு.
\end{defn}

\begin{explain}
தொகுத்தறிதல் வார்ப்புருவின் ஒவ்வொரு நிகழ்வும்~$\Struct{N}$ இல் மெய்யாகும்.
!!{formula}~$!A$ க்கு ஒரே ஒரு கட்டற்ற !!{variable}~$x$ மட்டுமே இருந்தால்
இதை மிக எளிதாகக் காணலாம். அப்போது $!A(x)$ ஆனது~$\Struct{N}$ இல்~$\Nat$
இன் உட்கணம்~$X_{!A}$ ஐ வரையறுக்கிறது. $s(x) = n$ ஆகும்போது
$\Sat{N}{!A(x)}[s]$ என அமையும் எல்லா~$n \in \Nat$ களின் கணமே~$X_{!A}$.
இதற்கு ஒத்த தொகுத்தறிதல் வார்ப்புரு நிகழ்வு
\[
((!A(\Obj 0) \land \lforall[x][(!A(x) \lif !A(x'))]) \lif 
  \lforall[x][!A(x)]).
\]
அதன் முன்னுறுப்பு~$\Struct{N}$ இல் மெய்யாக இருந்தால், $0 \in X_{!A}$;
மேலும் $n \in X_{!A}$ எனும்போதெல்லாம் $n+1$ உம் அதில் உள்ளது.
$0 \in X_{!A}$ என்பதால் $1 \in X_{!A}$ கிடைக்கிறது. $1 \in X_{!A}$
என்பதிலிருந்து $2 \in X_{!A}$ கிடைக்கிறது. இவ்வாறே தொடர்கிறது. ஆகவே
ஒவ்வொரு $n \in \Nat$ க்கும் $n \in X_{!A}$. ஆனால் இதன் பொருள்
$\lforall[x][!A(x)]$ ஆனது~$\Struct{N}$ இல் நிறைவுறுத்தப்படுகிறது என்பதே.
\end{explain}

$\Th{Q}$, $\Th{PA}$ இரண்டும் அடிகோளாக்கப்பட்ட கோட்பாடுகள். அவை எவ்வளவு
வலிமையுள்ளவை என்பதே பெரிய கேள்வி. எடுத்துக்காட்டாக, $\Lang L_A$ இல்
வெளிப்படுத்தக்கூடிய~$\Nat$ பற்றிய எல்லா மெய்மைகளையும் $\Th{PA}$ ஆல்
நிறுவ முடியுமா? குறிப்பாக, $\Lang L_A$ இல் கூறக்கூடிய எல்லாக் கேள்விகளையும்
$\Th{PA}$ இன் அடிகோள்கள் தீர்க்கின்றனவா?

இதை வேறொரு வகையில் கூறுவதானால்: $\Th{PA} = \Th{TA}$ ஆகுமா?
$\Th{TA}$ ஆனது~$\Nat$ பற்றிய எல்லா மெய்மைகளையும் வெளிப்படையாகவே
நிறுவுகிறது (அதாவது, உள்ளடக்குகிறது). மேலும், $!A$ ஆனது~$\Lang L_A$
இலுள்ள !!a{sentence} எனில் $\Sat{N}{!A}$ அல்லது $\Sat{N}{\lnot !A}$;
ஆகவே $\Th{TA} \Entails !A$ அல்லது $\Th{TA} \Entails \lnot !A$.
எனவே $\Lang L_A$ இல் கூறக்கூடிய எல்லாக் கேள்விகளையும் அது தீர்க்கிறது.
அத்தகைய கோட்பாட்டை \emph{!!{complete}} என்க.

\begin{defn}
ஒரு கோட்பாடு $\Gamma$ ஆனது, அதன் மொழியிலுள்ள ஒவ்வொரு
!!{sentence}~$!A$ க்கும் $\Gamma \Entails !A$ அல்லது $\Gamma \Entails
\lnot !A$ எனில், அப்போதும் அப்போது மட்டுமே \emph{!!{complete}}.
\end{defn}

\begin{explain}
நிறைவு தேற்றத்தின்படி $\Gamma \Entails !A$ அப்போதும் அப்போது மட்டுமே
$\Gamma \Proves !A$. ஆகவே, அதன் மொழியிலுள்ள ஒவ்வொரு
!!{sentence}~$!A$ க்கும் $\Gamma \Proves !A$ அல்லது $\Gamma \Proves
\lnot !A$ எனில், அப்போதும் அப்போது மட்டுமே $\Gamma$ நிறைவானது.
\end{explain}

இது மற்றொரு கேள்விக்கு வழிவகுக்கிறது: ஓர் !!a{sentence} ஆனது~$\Th{TA}$
இலோ, $\Th{PA}$ இலோ, $\Th{Q}$ இலாவது உள்ளதா என்று சோதிக்க நாம்
பயன்படுத்தக்கூடிய கணக்கீட்டுச் செயல்முறை உள்ளதா? ஒரு கணம் (எடுத்துக்காட்டாக,
!!{sentence}s இன் கணம்) எப்போது !!{decidable} என்று வரையறுப்பதன் மூலம்
இதைத் துல்லியமாக்கலாம்.

\begin{defn}
உள்ளீடு~$x$ க்கு, $x \in X$ எனில்~$1$ ஐயும் வேறெனில்~$0$ ஐயும் திருப்பித்
தரும் கணக்கீட்டுச் செயல்முறை இருந்தால், அப்போதும் அப்போது மட்டுமே கணம் $X$
\emph{!!{decidable}}.
\end{defn}

ஆகவே நமது கேள்வி: $\Th{TA}$ ($\Th{PA}$, $\Th{Q}$) !!{decidable} ஆகுமா?

இக்கேள்விகள் அனைத்துக்கும் விடை: இல்லை. இக்கோட்பாடுகளில் எதுவும்
தீர்மானிக்கத்தக்கதல்ல. ஆனால் இந்நிகழ்வு இக்குறிப்பிட்ட கோட்பாடுகளுக்கு மட்டும்
உரியதல்ல. உண்மையில், சில நிபந்தனைகளை நிறைவுசெய்யும் \emph{எந்தக்}
கோட்பாட்டுக்கும் அதே முடிவுகள் பொருந்தும். $\Th{Q}$, $\Th{PA}$ ஆகியவை
நிறைவுசெய்யும் இந்நிபந்தனைகளில் ஒன்று, அவை !!a{decidable} அடிகோள் கணத்தால்
அடிகோளாக்கப்படுகின்றன என்பதாகும்.

\begin{defn}
ஒரு கோட்பாடு !!a{decidable} அடிகோள் கணத்தால் அடிகோளாக்கப்பட்டால், அது
\emph{!!{axiomatizable}} எனப்படுகிறது.
\end{defn}

\begin{ex}
!!{sentence}s இன் ஒரு முடிவுறு கணத்தால் அடிகோளாக்கப்படும் எந்தக் கோட்பாடும்
!!{axiomatizable}; ஏனெனில் எந்த முடிவுறு கணமும் !!{decidable}. ஆகவே,
எடுத்துக்காட்டாக $\Th{Q}$ !!{axiomatizable}.

$\Th{PA}$ போன்ற வார்ப்புருவால் அடிகோளாக்கப்பட்ட கோட்பாடுகளும்
!!{axiomatizable}. $!B$ ஆனது~$\Th{PA}$ இன் அடிகோள்களுள் ஒன்றா என்று
சோதிக்க, அதாவது $!B$ ஆனது~$\Th{PA}$ இன் அடிகோள் எனில்
$\Char{X}(!B) = 1$, வேறெனில் $= 0$ என அமையும் $\Char{X}$ சார்பைக்
கணிக்க, பின்வருமாறு செய்யலாம். முதலில் $!B$ ஆனது~$\Th{Q}$ இன்
அடிகோள்களுள் ஒன்றா என்று சோதிக்கவும். அவ்வாறெனில் விடை ``ஆம்,'' மேலும்
$\Char{X}(!B) = 1$. அவ்வாறில்லையெனில், அது தொகுத்தறிதல் வார்ப்புருவின்
ஒரு நிகழ்வா என்று சோதிக்கவும். இதை முறைப்படி செய்யலாம். இந்நேர்வில்,
பின்வருமாறு செய்யலாம் என்பதே அதை மிக எளிதாகக் காட்டக்கூடும்: தொகுத்தறிதல்
வார்ப்புருவின் எந்த நிகழ்வும் பல அனைத்தளவையடைகளுடன் தொடங்கி, பின்னர் ஒரு
நிபந்தனை வாய்பாடாக அமையும் துணை-!!{formula} ஐக் கொண்டுள்ளது. அந்த நிபந்தனை
வாய்பாட்டின் பின்னுறுப்பு $\lforall[x][!A(x, y_1, \dots, y_n)]$; இங்கு
$x$, $y_1$, \dots, $y_n$ ஆகியவை~$!A$ இன் எல்லாக் கட்டற்ற மாறிகள்;
$!B$ இன் தொடக்க அளவையடைகள்~$y_1$, \dots,~$y_n$ மாறிகளைக் கட்டுகின்றன.
இந்த~$!A$ ஐப் பிரித்தெடுத்து, அதன் கட்டற்ற மாறிகள் முன்புற அனைத்தளவையடைகள்
மற்றும்~$\lforall[x]$ கட்டும் மாறிகளுடன் பொருந்துகின்றனவா என்று
சரிபார்த்ததும், நிபந்தனை வாய்பாட்டின் முன்னுறுப்பு பின்வருவதுடன் பொருந்துகிறதா
என்று சரிபார்ப்போம்:
\[
!A(\Obj 0, y_1, \dots, y_n) \land \lforall[x][(!A(x, y_1, \dots, y_n)
\lif !A(x', y_1, \dots, y_n))]
\]
மீண்டும், அது பொருந்தினால் $!B$ தொகுத்தறிதல் வார்ப்புருவின் ஒரு நிகழ்வு;
பொருந்தாவிட்டால் $!B$ அவ்வாறில்லை.
\end{ex}

இக்கேள்விக்கும்---எந்தெந்த கோட்பாடுகள் நிறைவானவை அல்லது தீர்மானிக்கத்தக்கவை
என்ற பொதுவான கேள்விக்கும்---விடையளிக்க, பின்வரும் வரையறையையும் கருதுவது
பயனுள்ளதாக இருக்கும். ஒரு கணம்~$X$ வெற்றாகவோ, !!a{surjective}
சார்பு~$f \colon \Nat \to X$ உள்ளதாகவோ இருந்தால், அப்போதும் அப்போது
மட்டுமே அக்கணம் !!{enumerable} என்பதை நினைவுகூர்க. அத்தகைய சார்பு~$X$ இன்
எண்ணிடல் எனப்படுகிறது.

\begin{defn}
ஒரு கணம்~$X$ வெற்றாகவோ, கணிக்கத்தக்க எண்ணிடலைக் கொண்டதாகவோ இருந்தால்,
அப்போதும் அப்போது மட்டுமே அது \emph{!!{computably enumerable}}
(சுருக்கமாக !!{c.e.}) எனப்படுகிறது.
\end{defn}

!!{axiomatizability} உடன், முழுமையின்மைத் தேற்றங்கள் பொருந்தும்
கோட்பாடுகளுக்கான மற்றொரு நிபந்தனை, கணிக்கத்தக்க சார்புகள் மற்றும்
!!{decidable} தொடர்புகள் பற்றிய அடிப்படை உண்மைகளை நிறுவும் அளவுக்கு அவை
வலிமையுள்ளதாக இருக்க வேண்டும் என்பதாகும். ``அடிப்படை உண்மைகள்'' என்பதால்,
கணிக்கத்தக்க சார்புகளின் ஒவ்வொரு உள்ளீட்டுக்கும் அவற்றின் மதிப்புகள் என்ன
என்பதை வெளிப்படுத்தும் !!{sentence}s ஐக் குறிக்கிறோம். ``போதிய
வலிமையுள்ளது'' என்பதால், குறித்த கோட்பாடுகள் இவ்வாக்கியங்களைத் தமது
தேற்றங்களுள் கொண்டுள்ளன என்று குறிக்கிறோம். எடுத்துக்காட்டாக, கணிக்கத்தக்க
சார்பின் முன்மாதிரியான கூட்டலைக் கருதுக. $2$, $3$ என்ற உள்ளீடுகளுக்கு~$+$
இன் மதிப்பு~$5$; அதாவது $2+3 = 5$. $2$, $3$ என்ற உள்ளீடுகளுக்கு~$+$ இன்
மதிப்பு~$5$ என்று வெளிப்படுத்தும் எண்கணித மொழி வாக்கியம்:
$(\num{2} + \num{3}) = \num{5}$. எடுத்துக்காட்டாக, இவ்வாக்கியத்தை
$\Th{Q}$ நிறுவுகிறது. இன்னும் பொதுவாக, ஒவ்வொரு கணிக்கத்தக்க சார்பு
$f(x_1, x_2)$ க்கும், $f(n_1, n_2) = m$ எனும்போதெல்லாம்
$\Th{Q} \Proves !A_f(\num{n_1}, \num{n_2}, \num{m})$ என அமையும்
!!a{formula}~$!A_f(x_1, x_2, y)$ ஆனது~$\Lang{L_A}$ இல் இருக்க வேண்டும்
என்று விரும்புகிறோம். இவ்வழியில், $n_1$, $n_2$ என்ற உள்ளீடுகளுக்கு~$f$
இன் மதிப்பு~$m$ என்று $\Th{Q}$ நிறுவுகிறது. உண்மையில், இன்னும் சற்றுக்
கூடுதலாக, $n_1$,~$n_2$ என்ற உள்ளீடுகளுக்கு வேறெந்த எண்ணும்~$f$ இன்
மதிப்பல்ல என்று அது நிறுவ வேண்டும். !!{decidable} தொடர்புகளுக்கும் இதுவே
பொருந்தும். பின்வரும் இரு வரையறைகள் இதைத் துல்லியமாக்குகின்றன.

\begin{defn}
!!{formula}~$!A(x_1, \dots, x_k, y)$ ஆனது, $f(n_1, \dots, n_k) = m$
எனும்போதெல்லாம் பின்வருவன நிலவினால், அப்போதும் அப்போது மட்டுமே
கோட்பாடு~$\Gamma$ வில் $f\colon \Nat^k \to \Nat$ என்ற சார்பை
\emph{!!{represents}} செய்கிறது:
\begin{enumerate}
\item $\Gamma \Proves !A(\num{n_1}, \dots, \num{n_k}, \num{m})$, மேலும்
\item $\Gamma \Proves \lforall[y](!A(\num{n_1}, \dots, \num{n_k},
y) \lif y = \num{m})$.
\end{enumerate}
\end{defn}

\begin{defn}
!!{formula}~$!A(x_1, \dots, x_k)$ ஆனது பின்வருவன நிலவினால், அப்போதும்
அப்போது மட்டுமே $R \subseteq \Nat^k$ என்ற தொடர்பை
\emph{!!{represents}} செய்கிறது:
\begin{enumerate}
\item $R(n_1, \dots, n_k)$ எனும்போதெல்லாம், $\Gamma \Proves
!A(\num{n_1}, \dots, \num{n_k})$, மேலும்
\item $R(n_1, \dots, n_k)$ அல்ல எனும்போதெல்லாம், $\Gamma \Proves \lnot
!A(\num{n_1}, \dots, \num{n_k})$.
\end{enumerate}
\end{defn}

ஒரு கோட்பாடு எல்லாக் கணிக்கத்தக்க சார்புகளையும் எல்லா !!{decidable}
தொடர்புகளையும் !!{represents} செய்தால், முழுமையின்மைத் தேற்றங்கள்
பொருந்தும் அளவுக்குப் ``போதிய வலிமையுள்ளது.'' $\Th{Q}$ உம் அதன்
விரிவாக்கங்களும் இந்நிபந்தனையை நிறைவுசெய்கின்றன. ஆனால் இதை நிறுவச் சிறிது
காலம் தேவைப்படும்---$\Th{Q}$ எவற்றை நிறுவ முடியும் என்பதைப் பற்றிய இது
எளிதல்லாத உண்மை; $\Th{Q}$ மிகச் சில அடிகோள்களை மட்டுமே கொண்டிருப்பதால்
அவற்றிலிருந்து இவ்வுண்மைகள் அனைத்தையும் நிறுவ வேண்டியுள்ளது. இருப்பினும், $\Th{Q}$
மிகவும் வலுக்குறைந்த கோட்பாடு. ஆகவே $\Th{Q}$ எல்லாக் கணிக்கத்தக்க
சார்புகளையும் பிரதிநிதித்துவப்படுத்துகிறது என்று நிறுவுவது கடினமாக இருந்தாலும்,
ஆர்வமூட்டும் பெரும்பாலான கோட்பாடுகள்~$\Th{Q}$ ஐவிட வலிமையானவை; அதாவது
$\Th{Q}$ நிறுவுவதைவிடக் கூடுதலாக நிறுவுகின்றன. $\Th{Q}$ ஒன்றை நிறுவினால்,
அதைவிட வலிமையான எந்தக் கோட்பாடும் அதை நிறுவும்; $\Th{Q}$ எல்லாக்
கணிக்கத்தக்க சார்புகளையும் பிரதிநிதித்துவப்படுத்துவதால், ஒவ்வொரு வலிமையான
கோட்பாடும் அவ்வாறே செய்கிறது. எனவே ஆர்வமூட்டும் பல கோட்பாடுகள்
முழுமையின்மைத் தேற்றங்களின் இந்நிபந்தனையை நிறைவுசெய்கின்றன. இவ்வாறு
இத்தேற்றங்கள் மிகப் பல கோட்பாடுகளுக்குப் பொருந்துகின்றன என்று காட்டுவதால்,
நமது கடினப் பணி பயனளிக்கும். ``கணிதம் முழுவதையும்'' முறைப்படுத்தும் நோக்குள்ள
எந்தக் கோட்பாடும் $\Th{Q}$ நிறுவும் அனைத்தையும் நிச்சயமாக நிறுவ வேண்டும்;
ஏனெனில் தொடக்கநிலைக் கணக்கீடுகளின் முடிவுகளையாவது அது பற்றிக்கொள்ள வேண்டும்.
ஆகவே ``கணிதம் முழுவதற்குமான'' கோட்பாடாக அமையும் எந்த வேட்புக் கோட்பாட்டுக்கும்
முழுமையின்மைத் தேற்றங்கள் பொருந்தும்.


\end{document}
